%% Elsevier elsarticle template derivative; LPPL 1.2 or later.
%% Copyright 2007, 2008, 2009 Elsevier Ltd.
%% Template derivative distributable under LPPL version 1.2 or later.
%% Software and original prose have separate licenses; see repository notices.
\documentclass[preprint,12pt,a4paper]{elsarticle}
\usepackage{amssymb}
\usepackage{amsmath}
\usepackage{array}
\usepackage{hyperref}
\hypersetup{hidelinks,pdftitle={OrchLang: Contract-Aware Resource Auditing for Bounded LLM Workflows}}
\journal{Science of Computer Programming}
\newcommand{\evCppTests}{106}
\newcommand{\evCostWorkflows}{23}
\newcommand{\evCostRuns}{115}
\newcommand{\evGenerated}{576}
\newcommand{\evAccepted}{24}
\newcommand{\evAcceptedFiniteCounterexamples}{0}
\newcommand{\evFiniteEqualRejections}{16}
\newcommand{\evTraceIdentityDisagreements}{0}
\newcommand{\evMutantsDetected}{4}
\newcommand{\evTokenizerEncodings}{24}
\newcommand{\evUnderestimates}{13}
\newcommand{\evConcatenations}{10}
\newcommand{\evNonadditive}{6}
\newcommand{\evCertificates}{24}
\newcommand{\evLeanChecks}{624}
\newcommand{\evSourceTampers}{48}
\newcommand{\evMalformedBatches}{4}
\newcommand{\evReplayCases}{12}
\newcommand{\evReplayPaths}{48}
\newcommand{\evExternalProjections}{11}
\newcommand{\evExternalRuns}{220}
\newcommand{\evExternalBoundViolations}{0}
\newcommand{\evExternalBaselineDisagreements}{0}
\newcommand{\evOriginalApplicationExecutions}{0}
\newcommand{\evSpentUsd}{0.00}

\newcommand{\evDataflowSourcePrograms}{256}
\newcommand{\evDataflowAccepted}{16}
\newcommand{\evDataflowExecutionContexts}{3072}
\newcommand{\evDataflowAcceptedFiniteCounterexamples}{0}
\newcommand{\evDataflowSourceReconstructionDisagreements}{0}
\newcommand{\evDataflowLeanNormalizationDisagreements}{0}
\newcommand{\evDataflowLeanReferenceDisagreements}{0}
\newcommand{\evDataflowLeanReplayDisagreements}{0}
\newcommand{\evDataflowCppAccepted}{8}
\newcommand{\evDataflowCppSecretArgumentCases}{192}
\newcommand{\evDataflowBoundaryCases}{8}
\newcommand{\evDataflowBoundaryPaths}{23}
\newcommand{\evDataflowSourceTampersRejected}{32}
\newcommand{\evDataflowAlteredPlansRejected}{16}
\newcommand{\evDataflowInvalidEvaluatorInputsRejected}{8}
\newcommand{\evDataflowCppTests}{106}
\newcommand{\evDataflowPythonTests}{47}
\newcommand{\evDataflowProofFiles}{3}

\newcommand{\revCppTests}{106}
\newcommand{\revPythonTests}{82}
\newcommand{\revProofFiles}{4}
\newcommand{\revPrecisionSourceCases}{80}
\newcommand{\revPrecisionSourceAccepted}{64}
\newcommand{\revPrecisionSourceBaseline}{8}
\newcommand{\revPrecisionCoreCases}{72}
\newcommand{\revPrecisionCoreAccepted}{64}
\newcommand{\revPrecisionCoreBaseline}{8}
\newcommand{\revPrecisionContexts}{3648}
\newcommand{\revAggregateCases}{256}
\newcommand{\revAggregateAccepted}{28}
\newcommand{\revAggregateBaseline}{16}
\newcommand{\revAggregateGains}{12}
\newcommand{\revAggregateContexts}{512}
\newcommand{\revAggregatePaths}{2048}
\newcommand{\revTokenizerTexts}{5344}
\newcommand{\revTokenizerConcats}{8996}
\newcommand{\revExactSupported}{40}
\newcommand{\revExactPaths}{128}
\newcommand{\revExactViolations}{0}

\begin{document}
\begin{frontmatter}
\title{OrchLang: Contract-Aware Resource Auditing for Bounded LLM Workflows}
\author[ci]{Nambi Rajan M}
\ead{nambirajan.m2024@vitstudent.ac.in}
\author[analytics]{Ranjithkumar S\corref{cor}}
\ead{ranjithkumar.s@vit.ac.in}
\cortext[cor]{Corresponding author.}
\address[ci]{Department of Computational Intelligence, SCOPE, Vellore Institute of Technology, India}
\address[analytics]{Department of Analytics, SCOPE, Vellore Institute of Technology, India}
\begin{abstract}
A bound on a language-model workflow's resource use does not establish that its
resource observations are independent of a secret. OrchLang provides separate
analyses for these obligations, with resolved contracts, complete request
expressions and executable Lean decision procedures. Proved request and
control-flow simplifications admit \revPrecisionSourceAccepted{} of
\revPrecisionSourceCases{} constructed source cases, compared with
\revPrecisionSourceBaseline{} for syntactic normalization. A separate aggregate
judgment admits permutations of two independent closed calls. Exact pinned
tokenizers expose a concatenation underbound and support complete-request
accounting. Evaluation retains the original 576-program and 256-pipeline
campaigns, new precision tests, and original-framework/local-model comparisons
of six workflow variants. These distinguish one direct source port from five
compositions retaining host-language control. The artifact connects conditional
core guarantees to inspectable source and execution contracts and retains the
generation and integration mismatches encountered.
\end{abstract}
\begin{keyword}
resource analysis \sep information flow \sep workflow languages \sep
verified decision procedures \sep reproducibility
\end{keyword}
\end{frontmatter}
\clearpage

\section*{Metadata}
\begin{table}[!ht]
\small
\begin{tabular}{|l|>{\raggedright\arraybackslash}p{.37\linewidth}|>{\raggedright\arraybackslash}p{.48\linewidth}|}
\hline
\textbf{Nr.} & \textbf{Code metadata description} & \textbf{Description} \\
\hline
C1 & Current code version & Scientific revision, 24 October 2026. The review package records the exact Git revision and source-file hashes. \\
\hline
C2 & Permanent link to code/repository used for this code version & \href{https://github.com/guyoverclocked/orchlang-compiler-design-project}{Project repository}; public release metadata for this local revision is pending. \\
\hline
C3 & Permanent link to Reproducible Capsule & Local source, evidence and manuscript archives for supervisor review; archival deposition pending. \\
\hline
C4 & Legal Code License & MIT software; LPPL template derivative; upstream notices retained. \\
\hline
C5 & Code versioning system used & Git \\
\hline
C6 & Software code languages, tools, and services used & C++17, Python, Lean 4.19.0; optional Tokenizers 0.21.4, LangGraph, Prompt Flow and local Ollama/Qwen. \\
\hline
C7 & Compilation requirements, operating environments and dependencies & CPU; C++17 and Python 3.10+. Recorded validation: Windows 11/GCC 16.1 and Ubuntu 24.04 under WSL2/GCC 13.3. Optional versions are pinned. \\
\hline
C8 & If available, link to developer documentation/manual & Repository README; \path{docs/REPRODUCIBILITY.md}, \path{docs/EXACT_ACCOUNTING.md}, and \path{docs/REAL_WORKFLOWS.md}. \\
\hline
C9 & Support email for questions & \texttt{ranjithkumar.s@vit.ac.in} \\
\hline
\end{tabular}
\caption{Review-package metadata. Submission administration is separate from the scientific evaluation.}
\end{table}

\section{Motivation and significance}
Consider a multi-tenant support service that uses a private customer attribute
to select an escalation workflow. A tenant-facing accounting service or
operations contractor sees per-model input usage, output usage and call counts,
but cannot read the attribute. Avoiding explicit disclosure in prompts does not
prevent the selected model or number of calls from revealing the branch. A
developer needs two answers: whether consumption is bounded, and whether the
specified resource observation changes with that secret.

Published attacks establish that encrypted language-model traffic can reveal
information through metadata. Weiss et al.\ study token-length sequences;
Whisper Leak studies packet sizes and timing~\cite{weiss,whisper}.
The accounting observer above is a deployment scenario, not an empirical
finding that these attacks use invoices. Our formal observation records
resource events; it excludes timing and network packets.

OrchLang is a small language and auditing artifact for bounded workflows.
Its central design choice is to retain complete constructed requests and
semantic contracts. Equal input lengths or maximum costs do not imply equal
response distributions. Later requests must retain the actual response strings
on which they depend.

The software provides (i) a bounded compiler with resolved contracts and
requests and diagnostic counterexamples; (ii) executable Lean decisions with
proved normalization and an aggregate-only commutation rule; (iii) pinned
complete-request accounting with explicit compositional assumptions; and
(iv) reproducible precision, accounting and original-framework execution
studies. These specialize established techniques to inspectable workflows.

\subsection{Relation to prior work}
Resource-aware noninterference combines information flow and resource typing;
RelCost bounds relative costs; probabilistic AARA supports expected costs;
and eRHL reasons about probabilistic program
relations~\cite{ngo,relcost,wang,erhl}. Stochastic calls therefore do not preclude
numeric specifications. Typed Resources supplies machine-checked scalar/vector
bounds for language-model workflows, while Costwright analyzes existing
framework code and provider-specific caps~\cite{typedresources,costwright}.
Token Budgets studies production overruns and affine budget
ownership~\cite{tokenbudgets}. OrchLang addresses a separate practical judgment:
invariance of resource observations when a secret changes, under declared
request-kernel assumptions.

CaMeL separates trusted control from untrusted data and tracks capabilities;
FIDES develops confidentiality/integrity policies for agent
planners~\cite{camel,fides}. CaMeL explicitly discusses side channels.
Hard-KLAIM proves typed trace safety at a declared interface; PC2 checks
obligations before operations and identifies model/world correspondence
conditions~\cite{hard,pc2}. OrchLang's contribution is a restricted executable
analysis, rather than a new general logic or an expressiveness separation.

Loom derives structural resource bounds and reports live-model
measurements~\cite{loom}. We compare with a competent sum/max/retry-multiply
baseline and add actual local generation. A source audit and two retained
probes of Costwright 0.2.15 did not reproduce OrchLang's historical identity
error; no external vulnerability finding is claimed.

\section{Software description}
\subsection{Architecture and judgments}
The C++ frontend supports bounded inputs, prompts, lexical models, conditionals,
bounded retries, outputs and tool sinks. It checks types, declared
confidentiality/integrity policies, structural resource bounds and relational
signatures. Sequential costs add, alternative bounds take maxima, and bounded
repetition multiplies a body bound. Resolved bindings distinguish shadowed
declarations. Canonical abstract contracts permit equivalent aliases while
retaining model revisions and caps.

A historical failure illustrates why identities matter. The base revision
\texttt{a49fcdb} contains an earlier checker that conflated models named
\texttt{m} despite different caps. An accepted secret branch consumed 2 versus
987 synthetic units at seed 1 while all 95 original tests passed. The unary
bound covered both executions; the relational claim failed. The repaired
checker rejects this witness. It motivates the bug class rather than serving
as the paper's principal contribution.

\begin{figure}[!ht]
\centering
\fbox{\begin{minipage}{.92\linewidth}
Source $\longrightarrow$ C++ policy/resource diagnostics and structural IR.\\[3pt]
Restricted source $\longrightarrow$ trusted Python reader
$\longrightarrow$ decoded Lean core $\longrightarrow$ verified resource decision.\\[3pt]
Complete requests $\longrightarrow$ pinned tokenizer
$\longrightarrow$ finite accounting or local-model admission.\\[3pt]
Pinned upstream workflows $\longleftrightarrow$ executable OrchLang kernels
and explicit host composition $\longrightarrow$ retained traces.
\end{minipage}}
\caption{Separate evidence paths. Source readers and execution adapters are
tested implementations, not mechanically verified C++ translations.}
\end{figure}

The ordinary \texttt{orchc check} judgment remains the conjunction of the C++
checks. Separate dataflow and aggregate commands answer explicitly named
resource questions. They do not waive confidentiality errors or authorize
disclosure. The historical \texttt{certify} command emits an analysis report.
Saved JSON plans are diagnostic records: checking them recomputes normalization.
The primary formal interface is a \emph{verified decision procedure}, without
a search-versus-check complexity claim.

Independent readers reconstruct closed and response-dependent source fragments.
The dataflow reader admits shared public text/Boolean inputs, secret Boolean
guards and actual response bindings, with lexical names aligned to slots.
Secret text, source retry, tools and reclassification are rejected; outputs
must be root literals. A separate deployment interpreter executes sequential
call kernels with real responses and dynamic text output. Its scope differs
from the proved relational fragment.

\subsection{Ordered resource decision}
Let $K(c,r)$ be a shared stateless finite kernel indexed by semantic contract
$c$ and complete request $r$. Outcomes contain response text, failure status
and weight. Calls update response slots; sequence composes stores and retains
billed event traces. Events record contract, UTF-8 input/output bytes and
failure status. Core retry restores the outer store between attempts, bills
failures and stops after an attempt without a failed call.

For public Boolean store $l$, secret Boolean store $h$ and text store $\sigma$,
Lean proves
\[
 N(P)=\operatorname{some}(Q)
 \ \Longrightarrow\
 \mathcal{D}_{K,l,h}(P,\sigma)
 =\mathcal{D}^{\mathrm{plan}}_{K,l}(Q,\sigma).
\]
This equality concerns complete weighted text-state transformations.
Accepted programs therefore have equal observations for any two secret
Boolean stores with shared public values. Probability interpretation
additionally requires normalized nonnegative finite kernels.

The new normalizer proves append associativity, empty-string identities,
adjacent-literal folding, sequence-unit elimination, and collapse of public
branches whose normalized arms agree. Secret arms can consequently be
accepted when their original syntax differs. Request rewrites preserve the
fully constructed string. The source reader already folds some template
literals, so source-level and core-expression gains are measured separately.

\subsection{Aggregate resource decision}
An aggregate observer sums input bytes, output bytes, calls and failures by
contract. Exactly two independent closed calls $A$ and $B$ can be exchanged
under this observer even though their ordered traces differ. Requests must be
closed, kernels stateless, and outcomes sampled independently. This rule
excludes response-dependent calls, output contents, retries and timing.

The separate Lean development uses natural outcome weights and a positive
common denominator per action. Finite nonnegative rational kernels admit such
a representation. For every Boolean predicate $T$ of the full aggregate
observation, it establishes
\[
 \Pr[T(\operatorname{obs}(A;B))]
 =\Pr[T(\operatorname{obs}(B;A))].
\]
The mechanized theorem equates histogram numerators and product denominators.
A finite-list permutation proof justifies transposing the Cartesian outcome
enumeration. The decoded decision checks that both secret arms contain the
same two full actions, possibly exchanged. It has a separate command and
observer; the ordered judgment remains unchanged.

An additional theorem preserves the joint distribution of usage and a
full-action-indexed multiset of actual response text/status pairs. A separate
finite execution mode can therefore retain results for a trusted host.
Source branch bindings still cannot escape their lexical branch; this is
not an end-to-end application or public-output confidentiality theorem.

\subsection{Exact tokenizer accounting}
The exact CLI takes a tokenizer JSON, its SHA-256 digest, a pinned library
version and an explicit special-token policy. It encodes complete requests
with padding and truncation disabled. Closed-source mode computes exact
inputs and conditional output-cap bounds directly from supported
\texttt{.orch} files. A declared mock model remains an abstract source contract.

Adding counts of prompt pieces is unsafe. Under the pinned plain GPT-2 policy,
\[
 t(\texttt{hello})+t(\texttt{world})=2
 <t(\texttt{helloworld})=3.
\]
Plain BERT has an analogous example with \texttt{a} and \texttt{world}.
Dynamic static bounds instead use separately enforced UTF-8 byte contracts
under the audited plain GPT-2 byte-BPE policy. Each initial symbol consumes a
byte, and merges cannot increase their number; thus
$t(s)\leq|\operatorname{UTF8}(s)|$ for that pinned configuration.
This implementation argument is separate from the Lean theorem.
BERT has no asserted byte rule: normalization can remove arbitrarily many
control bytes, so output token caps alone cannot bound response bytes.

The finite executor checks exact output caps and required byte caps for every
table outcome, accounts for billed failures and bounds retry. Runtime-only
mode reserves the complete request count plus output cap before dispatch
and can refuse a later request. C++ mock measurements retain their original
synthetic unit.

\section{Illustrative examples}
The permanent negative case uses \texttt{yes} and \texttt{no!}, which have equal
byte length but can induce one versus 100 output bytes under a valid kernel.
A downstream witness returns \texttt{aa} or \texttt{bb}; subsequent
\texttt{Next:aa} and \texttt{Next:bb} requests have different output
distributions. The dataflow decision retains this dependency and rejects
replacing it with a fixed request.

A positive source case makes the same call in both alternatives of a public
conditional inside one secret arm. The other secret arm calls directly.
Proved public-branch collapse admits these different syntax trees.
For the aggregate decision, one arm performs two independent closed document
checks and the other exchanges them. Aggregate acceptance preserves totals;
the ordered decision rejects the schedule difference. A response dependency
removes the commutation justification. Three executable source illustrations
accompany the aggregate checker.

\section{Impact and evaluation}
The evaluation asks whether the software gains precision under its specified
observers, accounts for concrete tokenizer behavior, and preserves behavior in
the source fragments it can express. New raw data are retained separately from
historical campaigns. Protocols, source hashes, commands, failed attempts,
dependency/model identities and generated tables accompany the artifact.
Precision sets are constructed stress tests, not samples of production
prevalence.

\subsection{Implementation and regression evidence}
Current validation runs \revCppTests{} C++ tests and \revPythonTests{} Python
tests; optional dependency skips are explicitly counted in the records.
All \revProofFiles{} Lean files compile under 4.19.0 without admitted lemmas
or project-specific axioms. The updated dataflow proof additionally uses
standard \texttt{Classical.choice}; the inventory retains
\texttt{propext} and \texttt{Quot.sound}. Cross-platform checks and sanitizer
commands are retained. Ubuntu and Windows CI configuration is included.

The closed-request rerun covers 576 programs and accepts 24.
No accepted case has a counterexample under the supplied kernels.

The original pipeline campaign covers 256 programs. It retains 16 accepts
across 3,072 contexts, with agreement between Lean,
the reference executor and replay. C++ accepts eight. The eight additional
resource-only accepts have IDs 0, 34, 68, 102, 136, 170, 204 and 238.
Each feeds a branch-local response into a later prompt; C++ conservatively
assigns the secret PC label and emits E230 in both arms. The supplement
tabulates every source and diagnostic. Resource equality does not override
this confidentiality policy.

Four historical compiled semantic mutants are detected. A competent structural
baseline agrees on 23 existing workflows and 115 mock executions. The earlier
11 cost projections and 220 mock runs remain historical shape checks and are
not counted as faithful application translations.

\subsection{Precision beyond syntactic identity}
\begin{table}[!ht]
\centering\small
\begin{tabular}{lrrrr}
\hline
Grammar/observer & Cases & Baseline & Revised & Gain \\
\hline
Ordered source & 80 & 8 & 64 & 56 \\
Ordered core & 72 & 8 & 64 & 56 \\
Aggregate closed & 256 & 16 & 28 & 12 \\
\hline
\end{tabular}

\caption{Constructed precision cases. The aggregate row changes the observer;
it does not preserve ordered traces.}
\label{tab:precision}
\end{table}
The ordered source campaign accepts \revPrecisionSourceAccepted{} of
\revPrecisionSourceCases{} cases versus \revPrecisionSourceBaseline{} for the
previous syntactic procedure. The separate core-expression campaign accepts
\revPrecisionCoreAccepted{} of \revPrecisionCoreCases{} versus
\revPrecisionCoreBaseline{}. Together they cover \revPrecisionContexts{}
exact execution contexts with zero accepted supplied-kernel counterexamples.
Separating negatives retain changed requests, contracts and dependencies.
The aggregate campaign accepts \revAggregateAccepted{} of
\revAggregateCases{} versus \revAggregateBaseline{} for ordered identity.
All \revAggregateGains{} gains have different ordered distributions under the
supplied kernels, as expected. Its \revAggregateContexts{} contexts and
\revAggregatePaths{} weighted paths have no accepted aggregate counterexample.
These counts establish implemented differences within disclosed grammars,
not typical application gains or superiority over general relational logics.

\subsection{Pinned-tokenizer evidence}
The corpus contains \revTokenizerTexts{} text-policy and
\revTokenizerConcats{} concatenation-policy observations. Among 2,249 pairs
per plain policy, joined counts exceed component sums in 36 GPT-2 and 75 BERT
cases. Targeted counterexamples refute additivity rather than estimate an error
rate. Of 64 finite program/policy/guard configurations,
\revExactSupported{} have supported static bounds; 24 correctly lack a
required dynamic byte rule. The \revExactPaths{} supported enumerated paths
have \revExactViolations{} bound violations. Four closed-source observations
and two CLI checks exercise actual \texttt{.orch} input.

\subsection{Original-framework and local-model execution}
The amended experiment compares 24 original/port pairs (48 completed executions; 0 failed), with 88 actual generations and no replay in either arm. Final outputs and complete request/response multisets agree in 22 and 22 pairs. There are 0 recorded deployment-bound violations and 0 tokenizer/server input-count mismatches. The six source variants represent five upstream patterns: one direct call-kernel port and five explicit host compositions. Table~\ref{tab:local} reports observed costs and conservative bounds.

Captured-response replay passes 24/24 comparisons and matches 44 complete requests with zero new generations. This tests port conformance under common actual outcomes. A post-hoc repeatability control uses 16 actual generations in 8 identical-request pairs. Responses differ in 1 of the 8 pairs. Earlier independent-generation differences remain retained, so a matching rerun is not a seeded-determinism guarantee.

\begin{table}[!ht]
\centering\small
\begin{tabular}{lrrrr}
\hline
Source variant & Pairs & Equal & Peak tokens & Bound range \\
\hline
Chain graph & 4 & 4 & 41 & 1608--1728 \\
Chain functional & 4 & 4 & 187 & 1608--1728 \\
Parallel & 4 & 4 & 200 & 1608--1728 \\
Routing & 4 & 4 & 125 & 1072--1152 \\
Basic (direct) & 4 & 2 & 75 & 536--576 \\
NER (hybrid) & 4 & 4 & 127 & 536--576 \\
\hline
\end{tabular}
\caption{Actual local generation: peak input-plus-generated-output usage across both arms and caps, versus guarded deployment bounds.}
\label{tab:local}
\end{table}


The deployment adapter explicitly serializes requests, hashes its tokenizer
and pins the cached \texttt{qwen2.5:3b} manifest. It requests CPU generation
with temperature zero and one seed. Complete input counts are compared with
server prompt counts; output usage measures generated tokens rather than
re-encoding decoded text. Input ceilings, output caps, context limits and
maximum calls govern admission. This bound is a deployment contract, not an
exact-token interpretation of the C++ byte/divisor estimate.

The corpus retains upstream prompts, response dependencies, predicates,
routing and postprocessing. One direct source port is distinguished from five
compositions retaining host predicates, structured routing, scheduling or list
processing. Source annotation and host-glue counts are reported; no
retrospective human adaptation-time estimate is invented. Original engines
and ports share the deployment contract. Independent generation,
captured-response replay, repeatability controls and task quality are
reported as distinct endpoints.

Both scored tasks expose model failures: NER splits the phone number into
three entities, and routing selects a joke for the poem request. Each task
passes four of eight execution-level exact checks across both arms. Agreement
between implementations therefore does not imply task success. Open-ended
generation is unscored; 44 of the 88 generations reach their output cap.

\begin{table}[!ht]
\centering\small
\begin{tabular}{lrrrr}
\hline
Source variant & Original & Kernels & Input/cap & Host case \\
\hline
Chain graph & 42 & 16 & 2/2 & 13 \\
Chain functional & 30 & 16 & 2/2 & 13 \\
Parallel graph & 43 & 21 & 3/3 & 9 \\
Routing graph & 71 & 14 & 2/2 & 14 \\
Basic & 106 & 7 & 1/1 & 2 \\
NER & 56 & 8 & 2/1 & 5 \\
\hline
\end{tabular}

\caption{Descriptive nonempty source lines (including comments), executable
kernel lines, source input/cap annotations, and host case lines. Shared drivers
and the kernel interpreter are counted separately in the supplement. These
counts do not measure human adaptation time or code reduction.}
\label{tab:annotations}
\end{table}

\subsection{Threats to validity}
Python readers, JSON handling, tokenizer libraries, execution adapters and the
platform remain trusted. Lean proves decoded-core transformations, not C++
soundness, reader correctness or provider conformance. Probability
interpretation requires valid finite kernels; the relation assumes stateless
request-indexed behavior. Local caching and repeatability are examined
empirically rather than assumed to satisfy that theorem.

Precision grammars intentionally contain the supported equivalences.
Finite checks cannot replace soundness proofs or establish completeness.
The local corpus contains five distinct tutorial patterns, two selected inputs
per variant, one seed and one model. Hybrid compositions expose source support
gaps and are not whole-program translations. Resource compliance, trace
conformance and task quality answer different questions. Admission can refuse
a request before completion; failures require an explicit accounting contract.
The final local run follows disclosed infrastructure and translation repairs;
it is an amended development study, not untouched preregistration.
Timing, network metadata, secret text, mutable providers and production
deployment remain outside the demonstrated guarantee.

\section{Conclusions}
OrchLang makes resource contracts inspectable where source structure,
constructed requests, kernels and execution meet. Its verified decisions
admit specified semantic rewrites and independent-call aggregate permutations
beyond syntactic identity. Complete-request tokenization removes an unsafe
additive abstraction. Original-framework runs and separately labeled replay
checks expose both execution support and remaining host/model assumptions.
The artifact supplies executable decisions and counterexamples for developers
auditing these bounded workflow obligations.

\section*{Data, software and research-assistance statements}
The local review package includes source, protocols, raw data, proofs, claim
ledger and table scripts. Credentials, weights and private literature are
excluded. Public archival metadata and final author declarations remain
separate administrative steps. No paid inference API was used; API expenditure
is USD 0.

OpenAI Codex assisted implementation, proofs, experiments, literature retrieval
and writing. Git history also records earlier Anthropic Claude assistance
under the attribution \texttt{Claude Opus 5}. These systems are research tools,
not authors. Automated checks do not constitute completed human verification
or final author approval.

\begin{thebibliography}{20}
\bibitem{ngo} V. C. Ngo, M. Dehesa-Azuara, M. Fredrikson, J. Hoffmann,
Verifying and Synthesizing Constant-Resource Implementations with Types,
IEEE Symposium on Security and Privacy (2017).
\url{https://www.ieee-security.org/TC/SP2017/papers/408.pdf}.
\bibitem{relcost} E. \c{C}i\c{c}ek, G. Barthe, M. Gaboardi, D. Garg, J. Hoffmann,
Relational Cost Analysis, POPL (2017). \url{https://doi.org/10.1145/3009837.3009858}.
\bibitem{wang} D. Wang, D. M. Kahn, J. Hoffmann, Raising Expectations:
Automating Expected Cost Analysis with Types, Proc. ACM Program. Lang. 4,
ICFP, Article 110 (2020). \url{https://doi.org/10.1145/3408992}.
\bibitem{erhl} M. Avanzini, G. Barthe, D. Davoli, B. Gr\'egoire,
A quantitative probabilistic relational Hoare logic, arXiv:2407.17127v2 (2025).
\url{https://arxiv.org/abs/2407.17127v2}.
\bibitem{hard} M. Aldinucci, A. Bracciali, I. Colonnelli, D. Medi\'c,
A. Mulone, L. Roversi, M. E. Santimaria, Hard-KLAIM: Hardening MCP-Enabled
Agentic Workflows with a Locality-Aware Type-and-Effect System for Trace Safety,
Journal of Logical and Algebraic Methods in Programming, pre-proof (2026).
\url{https://doi.org/10.1016/j.jlamp.2026.101191}.
\bibitem{pc2} C. Ciliberto, A. Rudi, E. Tassi, Proof-Carrying Contracts:
Commitments AI Agents Cannot Break, arXiv:2610.28484v1 (2026).
\url{https://arxiv.org/abs/2610.28484v1}.
\bibitem{loom} K. Dickerson, Statically Bounding Resource Accumulation in LLM
Agent Workflows: The Cost Case, author-hosted report, 18 May 2026.
\url{https://loom.technology/research/bounded-llm-cost/}.
\bibitem{promptflow} Microsoft, Prompt Flow, standard examples, revision
\nolinkurl{33a68c9d3cf349d7fdfad1f41969af27543c0406}, accessed 23 October 2026.
\url{https://github.com/microsoft/promptflow}.
\bibitem{langgraph} LangChain, LangGraph workflows and agents documentation,
revision \nolinkurl{ecda7f522a6ffd5faaabb843b581fe57cd5eb782}, accessed 23 October 2026.
\url{https://github.com/langchain-ai/docs}.
\bibitem{costwright} H. Inverso, Costwright 0.2.15, inspected revision
\nolinkurl{2531c37cc2aa06472a2628bccfb527f185f6b4c0} (2026).
\url{https://github.com/eleata/costwright}.
\bibitem{typedresources} H. Inverso, A Potential-Based Calculus for Resource
Bounds of LLM-Agent Workflows, software and preprint (2026).
\url{https://doi.org/10.5281/zenodo.20661092}.
\bibitem{tokenbudgets} S. Khan, Token Budgets: An Empirical Catalog of 63
LLM-Agent Budget-Overrun Incidents, with an Affine-Typed Rust Mitigation as a
Case Study, arXiv:2606.04056 (2026). \url{https://arxiv.org/abs/2606.04056}.
\bibitem{camel} E. Debenedetti et al., Defeating Prompt Injections by Design,
arXiv:2503.18813 (2025). \url{https://arxiv.org/abs/2503.18813}.
\bibitem{fides} M. Costa et al., Securing AI Agents with Information-Flow
Control, arXiv:2505.23643 (2025). \url{https://arxiv.org/abs/2505.23643}.
\bibitem{weiss} R. Weiss, D. Ayzenshteyn, G. Amit, Y. Mirsky,
What Was Your Prompt? A Remote Keylogging Attack on AI Assistants,
USENIX Security (2024), 3367--3384.
\url{https://www.usenix.org/conference/usenixsecurity24/presentation/weiss}.
\bibitem{whisper} G. McDonald, J. Bar Or, Whisper Leak: a side-channel attack
on Large Language Models, Microsoft Research (November 2025).
\href{https://www.microsoft.com/en-us/research/publication/whisper-leak-a-side-channel-attack-on-large-language-models/}{Microsoft Research publication record}.

\end{thebibliography}

\end{document}
